\section{Proposed explainable algorithm}\label{sec:explainablealgorithm}
% No se comenta por qué se propone el nuevo algoritmo.
% Ni que lo hace diferente a los que aparecen en el trabajo relacionado.
% Además debería darse una explicación intuitiva de la idea del algoritmo.

{\color{blue}
     Current approaches for predicting credit risk achieve high levels of accuracy by applying techniques such as class balancing and feature selection to improve predictive performance. Recently, research has increasingly focused on explaining these predictions. Techniques such as SHAP and LIME attribute importance to individual features, but they do not indicate how these features interact in the prediction. This limitation highlights the need for algorithms that assess the contribution of each feature and explain predictions based on the interaction between these features.\\
    
        Therefore, we propose a pattern-based algorithm for predicting credit risk. Unlike other state-of-the-art algorithms, this algorithm allows predictions to be explained by identifying combinations of features (patterns). In this context, a pattern is a combination of appears frequently in a dataset \cite{gutierrez2015mining, jia2020exploiting}. \\
    
        The proposed algorithm consists of two stages. In the first stage, a state-of-the-art predictor is trained for credit risk prediction, and within each class (defaulters and non-defaulters), instances are clustered into groups to extract and filter patterns. The second stage focuses on building an explanation: when a new instance arrives, the algorithm predicts its class and then identifies the most similar group within that predicted class to generate an explanation based on the patterns of that group.\\
    
        In the initial stage, the proposed algorithm first employs an oversampling method specifically designed for mixed-data types to balance the dataset. Then, a state-of-the-art predictor tailored to handle such data is trained on the balanced dataset. This approach aims to mitigate the bias toward the majority class. The second process groups instances within each class into clusters to identify patterns within each cluster. A filtering procedure to eliminate redundancies, retaining only the most representative and frequent combinations of features that define each group.\\

        The algorithm's second stage is dedicated to prediction and explainability. On receiving a new instance, the algorithm first predicts its class (defaulter or non-defaulter). It then identifies the most similar cluster within that predicted class. The final explanation is constructed from the most frequent and representative patterns of this matching cluster's training data and covers the new instance. A pattern is said to cover the instance if the combination of feature values it represents matches the instance’s characteristics. This approach explains the prediction through combinations of features that frequently appear in similar training instances.\\
    
        In summary, our proposed algorithm integrates a prediction stage that allows the use of the best-performing state-of-the-art credit risk predictor with a pattern-based explanation stage, which allows explaining the results of any predictor. This algorithm seeks to provide high predictive accuracy with explainability. Because it can incorporate different predictors, the algorithm remains flexible and adaptable to future advances in credit risk prediction.\\
}

\subsection{First stage}
{\color{blue}
	The first stage of the proposed algorithm consists of training a predictor for credit risk prediction and mining contrast patterns that will later be used to explain its predictions. Since, as reported in the related work section, explanation methods based on weights or importance scores of individual features have limitations in representing interactions among multiple features, the proposed approach will explain predictions through contrast patterns that can involve multiple features. \newline
}

For credit risk prediction, we propose using the FT4cip predictor \cite{canete2022ft4cip}, a decision tree based predictor capable of handling mixed and imbalanced datasets, which has reported strong performance on this type of data, making it suitable for credit risk prediction. Although FT4cip is designed to handle imbalanced datasets, the experiments conducted in this work indicated that FT4cip achieved better predictive performance when trained using a previously balanced training set. Thus, following the results reported in \cite{melchor2025oversampling}, we propose applying SMOTENC \cite{chawla2002smote, mukherjee2021smote} to it to balance the training set before predictor training.\newline

To construct explanations, we propose mining contrast patterns from similar instances within each class. However, instances belonging to a same class may be from different subtypes  characterized by different feature-value combinations. Therefore, we first build clusters of similar instances within each class and subsequently mine contrast patterns from each cluster, obtaining contrast patterns that describe the feature-value combinations observed among similar instances into the cluster that can later be used to explain credit risk predictions for new instances. To generate these clusters, we propose using the Kmodes algorithm \cite{huang1997clustering, huang1998extensions}, as it is specifically designed for mixed data and avoids the limitations of distance-based methods when handling non-numeric features. To determine the number of clusters, following \cite{dinh2019estimating, kuswardana2025comparison}, we propose using the Silhouette coefficient, which, according to these works, allows obtaining well-defined clusters. {\color{blue}By mining contrast patterns from each cluster rather than from the entire class, the resulting contrast patterns represent feature-value combinations frequently observed among similar instances within the cluster. Since these patterns are mined from clusters of similar instances and typically cover a substantial portion of the cluster, they provide a description of the characteristics shared by the instances in that cluster. Consequently, these patterns can later be used to explain credit risk predictions for new instances by identifying feature-value combinations shared between the new instance and similar instances in a cluster.}\newline 

For mining contrast patterns, we propose using CPMining, introduced in the PBC4cip algorithm \cite{loyola2017pbc4cip}, as a supervised contrast pattern miner that explicitly exploits class differences to identify them. CPMining is applied independently to each cluster by contrasting the cluster's instances with all instances of the opposite class. This design allows the mining of contrast patterns that characterize instances within each cluster while distinguishing them from the opposite class, instead of generating contrast patterns for the entire class, whose objective is to cover as many instances as possible. {\color{blue} As a result, the generated contrast patterns are cluster-specific and describe feature-value combinations observed among similar instances within each cluster. When a pattern covering a new instance is selected as an explanation, its conditions correspond to those previously observed combinations. Furthermore, because both FT4cip and contrast patterns rely on conjunctions of $feature\#value$, the resulting explanation follows the same based on rules representation used by the predictor.}\newline

Due to the large number of contrast patterns generated during mining, the filtering used in CPMining is applied to reduce redundant patterns. This filtering removes more specific patterns when a more general pattern covers an instances and when redundant conditions are present within the same pattern. Additionally, unlike the filtering used in CPMining, the proposed approach removes contrast patterns whose support in the opposite class is non-zero. {\color{blue}As a result, the final set contains only contrast patterns mined from similar instances within the cluster and not observed in the opposite class. This allows explanations for new instances to be constructed using only patterns that cover the instance and represent feature-value combinations exclusively associated with similar instances in the cluster and predicted class.}\newline

{\color{blue}			
	\subsection{Second stage}
	Given a new instance $x$, the prediction stage determines its credit risk class (defaulter or non-defaulter) using the trained FT4cip predictor. Once the predicted class $c$ has been determined, the explanation stage identifies the cluster associated with class c that is most similar to the new instance and uses the contrast patterns mined from that cluster to construct the explanation.\newline
	
	To identify the cluster most similar to a new instance $x$, we define a contrast pattern-based similarity score between $x$ and each cluster $C_{c,j}$ associated with the predicted class $c$. Instead of relying on traditional similarity measures for mixed data, the proposed approach evaluates similarity using the contrast patterns extracted from each cluster, since these contrast patterns were previously mined from similar instances within each cluster and describe combinations of $feature\#value$  observed among those instances and associated with the predicted class. Consequently, similarity is computed from the support of the contrast patterns extracted from a cluster that cover the new instance, relative to the total support of the patterns extracted from that cluster. The similarity score is defined in Equation~\ref{eq:cluster_similarity}.\newline
	
	\begin{equation}
		Sim(x, C_{c,j})
		=
		\frac{
			\sum\limits_{p \in \mathcal{P}_{c,j}(x)} supp(p)
		}{
			\sum\limits_{p \in \mathcal{P}_{c,j}} supp(p)
		},
		\label{eq:cluster_similarity}
	\end{equation}
	This similarity measure increases as the support accumulated by the contrast patterns covering the new instance increases. Since patterns with higher support are observed in a larger number of cluster instances, they contribute more to the similarity score. Furthermore, dividing by the total support of the patterns extracted from the cluster allows proportional comparisons across clusters. Consequently, higher similarity values indicate a greater correspondence between the new instance and the feature-value combinations represented by the cluster's contrast patterns.\newline
	
	The cluster with the highest similarity score is selected as the cluster most similar to the new instance $x$. Once the most similar cluster $C_{c,j^*}$ has been identified, the explanation is constructed using the contrast patterns extracted from that cluster. Let $\mathcal{P}_{c,j^*}(x)$ denote the set of contrast patterns from $C_{c,j^*}$ that cover instance $x$. From this subset, the algorithm selects the contrast pattern $p$ with the highest support, as defined in Equation~\ref{eq:pattern_selection}. Support is used as the selection criterion because patterns with higher support correspond to combinations of conditions observed in a larger number of similar instances within the selected cluster. Therefore, the resulting pattern provides an explanation based on conditions that are shared by multiple instances similar to the one being explained.
	
	\begin{equation}
		p^* = \arg\max_{p \in \mathcal{P}_{c,j^*}(x)} supp(p),
		\label{eq:pattern_selection}
	\end{equation}
	
	The resulting pattern constitutes the final explanation of the prediction and is expressed as a conjunction of $feature\#value$ conditions. This explanation identifies the conditions satisfied by the new instance that are also observed among the instances of the most similar cluster. Since the selected cluster contains the instances most similar to the new instance within the predicted class, the explanation is constructed from the characteristics shared with those instances. If multiple patterns achieve the same maximum support, the pattern containing the largest number of $feature\#value$ conditions is selected. This criterion provides a more detailed explanation by incorporating additional conditions associated with the instances covered by the pattern. Furthermore, because both FT4cip and contrast patterns are expressed as conjunctions of $feature\#value$ conditions, the resulting explanation follows the same rule-based representation used by the predictor.
	
}